\chapter{项目七：插件式命令行工具}

\section{项目目标}

本项目实现一个插件式命令行工具骨架。它支持多个子命令，例如 \code{wordfreq}、\code{fileindex}、\code{configcheck}。每个子命令有自己的参数解析和执行逻辑，主程序只负责分发。

目标用法：

\begin{lstlisting}[language=bash,caption={插件式 CLI 用法}]
tool wordfreq --input article.txt --top 20
tool fileindex --root ./src --top 10
tool configcheck --path server.conf
\end{lstlisting}

这个项目训练大型代码组织：接口、注册表、命令边界、错误码和测试替身。

\section{命令接口}

每个命令都实现同一个接口：

\begin{lstlisting}[language=C++,caption={命令接口}]
struct CommandContext {
    std::ostream& output;
    std::ostream& error;
};

class Command {
public:
    virtual ~Command() = default;
    [[nodiscard]] virtual std::string_view name() const noexcept = 0;
    [[nodiscard]] virtual int run(std::span<const std::string_view> args,
                                  CommandContext context) const = 0;
};
\end{lstlisting}

为什么 \code{CommandContext} 里放输出流？这样测试命令时可以传 \code{std::ostringstream}，不用真的写终端。接口从一开始就为测试留出边界。

\section{命令注册表}

\begin{lstlisting}[language=C++,caption={命令注册表}]
class CommandRegistry {
public:
    void add(std::unique_ptr<Command> command)
    {
        const std::string key{command->name()};
        this->commands_[key] = std::move(command);
    }

    [[nodiscard]] const Command* find(std::string_view name) const
    {
        const auto found = this->commands_.find(std::string{name});
        if (found == this->commands_.end()) {
            return nullptr;
        }
        return found->second.get();
    }

private:
    std::unordered_map<std::string, std::unique_ptr<Command>> commands_;
};
\end{lstlisting}

注册表拥有命令对象，所以用 \code{std::unique_ptr<Command>}。查找返回裸指针，只表示观察，不转移所有权。这个接口把拥有和观察区分开了。

\section{主分发逻辑}

\begin{lstlisting}[language=C++,caption={命令分发}]
int dispatch(const CommandRegistry& registry,
             std::span<const std::string_view> args,
             CommandContext context)
{
    if (args.empty()) {
        context.error << "missing command\n";
        return 2;
    }

    const Command* command = registry.find(args.front());
    if (command == nullptr) {
        context.error << "unknown command: " << args.front() << '\n';
        return 2;
    }

    return command->run(args.subspan(1), context);
}
\end{lstlisting}

分发器不关心 \code{wordfreq} 怎么解析参数，也不关心 \code{fileindex} 怎么扫描文件。这就是边界。新增命令时，主分发逻辑不需要修改。

\section{实现一个 configcheck 命令}

\begin{lstlisting}[language=C++,caption={配置检查命令}]
class ConfigCheckCommand final : public Command {
public:
    [[nodiscard]] std::string_view name() const noexcept override
    {
        return "configcheck";
    }

    [[nodiscard]] int run(std::span<const std::string_view> args,
                          CommandContext context) const override
    {
        if (args.size() != 2 || args[0] != "--path") {
            context.error << "usage: configcheck --path FILE\n";
            return 2;
        }

        const ConfigResult result = load_config_file(std::filesystem::path{args[1]});
        if (result.error.has_value()) {
            context.error << result.error->message << '\n';
            return 1;
        }
        context.output << "ok\n";
        return 0;
    }
};
\end{lstlisting}

命令内部可以调用配置加载器项目。这样前面的项目不是孤立练习，而是能组合进更大的工具。

\section{注册命令}

\begin{lstlisting}[language=C++,caption={组装命令}]
CommandRegistry build_registry()
{
    CommandRegistry registry;
    registry.add(std::make_unique<ConfigCheckCommand>());
    registry.add(std::make_unique<WordFreqCommand>());
    registry.add(std::make_unique<FileIndexCommand>());
    return registry;
}
\end{lstlisting}

这里的 \code{build_registry} 是装配层。它知道具体命令类型，业务分发器不知道。装配层集中依赖具体实现，是大型项目常见组织方式。

\section{测试分发器}

可以写一个假命令：

\begin{lstlisting}[language=C++,caption={测试用假命令}]
class FakeCommand final : public Command {
public:
    [[nodiscard]] std::string_view name() const noexcept override
    {
        return "fake";
    }

    [[nodiscard]] int run(std::span<const std::string_view> args,
                          CommandContext context) const override
    {
        context.output << "args=" << args.size() << '\n';
        return 0;
    }
};
\end{lstlisting}

测试分发器时，不需要真实扫描文件或读取配置。只验证：缺命令报错，未知命令报错，已知命令收到剩余参数。这就是依赖抽象的测试价值。

\section{插件化的下一步}

本章的“插件式”还是编译期注册。真正动态插件还涉及动态库加载、ABI 稳定、版本协商和安全边界。不要一开始就做动态插件。先把命令接口、注册表和测试边界做好，已经能支撑大多数内部工具。

\begin{exercisebox}
给 CLI 增加 \code{help} 命令。要求每个命令提供 \code{usage()}，\code{help} 能列出所有命令。思考注册表是否应该保留命令顺序；如果需要稳定顺序，应该增加什么数据结构？
\end{exercisebox}

\section{项目从目录结构开始}

插件式 CLI 如果只停留在一个文件里，很快会变成大杂烩。一个可维护的最小目录可以这样安排：

\begin{lstlisting}[caption={命令行工具目录结构}]
src/
  main.cpp
  cli/
    command.hpp
    command_registry.hpp
    dispatch.cpp
  commands/
    config_check_command.cpp
    word_freq_command.cpp
    file_index_command.cpp
  app/
    build_registry.cpp
tests/
  cli/
    dispatch_tests.cpp
  commands/
    config_check_command_tests.cpp
\end{lstlisting}

\code{cli} 目录只放抽象接口和分发逻辑；\code{commands} 放具体命令；\code{app} 是装配层，知道要注册哪些命令；\code{main.cpp} 只连接系统入口。目录结构不是形式，它让依赖方向清楚：核心分发不依赖具体命令，具体命令依赖各自项目能力，装配层把它们拼起来。

\section{注册表还需要顺序}

前面的 \code{CommandRegistry} 用哈希表查找命令，适合分发。但 \code{help} 命令需要稳定列出命令。如果直接遍历 \code{unordered_map}，输出顺序不稳定。可以同时维护一个顺序数组：

\begin{lstlisting}[language=C++,caption={同时支持查找和稳定顺序}]
class CommandRegistry {
public:
    void add(std::unique_ptr<Command> command)
    {
        const std::string key{command->name()};
        this->ordered_.push_back(command.get());
        this->commands_[key] = std::move(command);
    }

    [[nodiscard]] std::span<Command* const> ordered() const noexcept
    {
        return this->ordered_;
    }

private:
    std::unordered_map<std::string, std::unique_ptr<Command>> commands_;
    std::vector<Command*> ordered_;
};
\end{lstlisting}

这里 \code{commands_} 拥有对象，\code{ordered_} 只保存观察指针。这个设计成立的前提是：命令对象由 \code{unique_ptr} 管理，移动 \code{unique_ptr} 本身不会移动被拥有的命令对象；注册表析构时，两个成员一起销毁，外部不能长期保存 \code{ordered} 返回的指针。如果你觉得这个生命周期证明太绕，可以改成 \code{std::vector<std::string>} 保存命令名，用名字再查找，牺牲一点效率换更简单的所有权。

\section{usage 是接口的一部分}

为了支持 \code{help}，命令接口应扩展：

\begin{lstlisting}[language=C++,caption={命令提供用法说明}]
class Command {
public:
    virtual ~Command() = default;
    [[nodiscard]] virtual std::string_view name() const noexcept = 0;
    [[nodiscard]] virtual std::string_view usage() const noexcept = 0;
    [[nodiscard]] virtual int run(std::span<const std::string_view> args,
                                  CommandContext context) const = 0;
};
\end{lstlisting}

然后 \code{HelpCommand} 可以依赖注册表的只读视图：

\begin{lstlisting}[language=C++,caption={help 命令列出所有命令}]
class HelpCommand final : public Command {
public:
    explicit HelpCommand(const CommandRegistry& registry)
        : registry_{registry}
    {
    }

    [[nodiscard]] std::string_view name() const noexcept override
    {
        return "help";
    }

    [[nodiscard]] std::string_view usage() const noexcept override
    {
        return "help";
    }

    [[nodiscard]] int run(std::span<const std::string_view>,
                          CommandContext context) const override
    {
        for (const Command* command : this->registry_.ordered()) {
            context.output << command->usage() << '\n';
        }
        return 0;
    }

private:
    const CommandRegistry& registry_;
};
\end{lstlisting}

这个版本引入了一个新生命周期要求：\code{HelpCommand} 引用的注册表必须比它活得久。如果 \code{HelpCommand} 本身也被放进注册表，就会形成自引用结构，设计要更谨慎。更简单的方案是让 \code{help} 不是普通命令，而是分发器内置逻辑。工程设计没有唯一答案，关键是把所有权和生命周期写清楚。

\section{错误码合同}

命令行工具的返回码也要统一：

\begin{longtable}{p{0.18\linewidth}p{0.62\linewidth}}
\toprule
返回码 & 含义 \\
\midrule
\code{0} & 成功 \\
\code{1} & 命令执行失败，例如文件不存在、配置非法 \\
\code{2} & 用法错误，例如未知命令、参数缺失、参数格式错误 \\
\bottomrule
\end{longtable}

有了这个合同，测试就能写得具体：

\begin{lstlisting}[language=C++,caption={分发器错误码测试}]
std::ostringstream output;
std::ostringstream error;
const int code = dispatch(registry, {}, CommandContext{output, error});
assert(code == 2);
assert(error.str().find("missing command") != std::string::npos);
\end{lstlisting}

返回码不是小事。脚本、CI 和其他程序都会依赖它判断工具是否成功。

\section{从编译期插件到动态插件}

本章的插件是编译期插件：所有命令在编译时链接进程序，只是在运行时通过注册表分发。它已经足够训练抽象边界。动态插件则是另一个项目：要考虑动态库加载、符号导出、ABI 稳定、版本协商、崩溃隔离和安全策略。不要把“可扩展”一上来理解成“必须加载动态库”。大多数内部 CLI 先做到编译期可扩展、测试清楚、错误码稳定，就已经很有价值。

\section{项目验收清单}

完成这个项目时，至少验收：

\begin{itemize}
  \item 缺命令返回 \code{2}，错误输出可读。
  \item 未知命令返回 \code{2}，包含原始命令名。
  \item 已知命令只收到自己的剩余参数。
  \item 命令可以用假输出流测试，不依赖真实终端。
  \item \code{help} 输出顺序稳定。
  \item 新增命令只需要改装配层和测试，不需要改分发器。
\end{itemize}

如果这些都成立，这个 CLI 骨架就不只是“能跑”，而是能承受后续项目继续接入。
